\subsection{Related Work}
\label{sec:related-work}
Our project intersects with four primary literatures: the techno-economic analysis literature on the impacts of energy efficiency on power consumption, the literature on power purchase agreements and their usage as a tool to reduce corporate carbon footprints, the literature on energy efficiency in AI, and the literature on power demand and power quality in the face of increasing load.

A substantial body of research has examined the techno-economic drivers of energy demand from data centers and AI workloads. Recent assessments by \cite{RN6} indicate that U.S. data center electricity consumption continues to rise, even as efficiency gains partially offset this growth. Complementary analyses from international organizations such as recent work by \cite{RN1} and industry-focused reports such as those from \cite{RN3} highlight the increasing share of artificial intelligence and machine learning workloads in total data center demand. These reports project that the rapid adoption of AI services will represent one of the dominant forces shaping future electricity requirements. At the micro-level, studies by \cite{RN9} and \cite{RN4} examine the energy performance of different AI accelerators and identify workload management opportunities for reducing power consumption. Together, these techno-economic assessments provide the empirical foundation for understanding how large-scale computing alters both system-level demand and the opportunities for efficiency improvements.

% In parallel, a growing literature has evaluated the role of corporate power purchase agreements (PPAs) in advancing renewable energy adoption. \cite{RN10, RN11} demonstrate that PPAs not only provide firms with cost stability but also have the potential to stimulate renewable capacity expansion. Their policy-oriented perspective emphasizes that the environmental effectiveness of PPAs depends critically on contract structure and market design. Extending this corporate decarbonization perspective, \cite{sirin2024energy} show that firms that adopt credible, capability-based energy transition strategies are rewarded in financial markets, underscoring the link between procurement choices and firm performance. Taken together, this literature positions PPAs as both a corporate strategy for risk management and a broader mechanism with systemic implications for renewable investment and decarbonization.

The environmental consequences of AI have received considerable scholarly attention. Early analyses showed that that large-scale neural network training would generate escalating carbon emissions across the AI life cycle~\cite{strubell19,schwartz2020green,RN7, RN5} and result in significant environmental and financial costs. Subsequent works~\cite{RN2, RN8} argue that efficiency improvements in hardware, software, and energy sourcing can cause the carbon footprint of AI training to plateau and ultimately decline. However, recent holistic assessments such as \cite{morrison2025holistically} show that large AI models continue to have considerable impact on carbon emission (and water usage), but conclude that the full extent of the environmental impact of AI is hard to quantify. %At the same time, research such as \cite{kim2025cost} emphasizes that inference-related costs—particularly under test-time scaling—are likely to become an increasingly important component of AI’s environmental burden. Broader assessments, including \cite{RN7, RN5}, employ life-cycle assessment (LCA) frameworks to highlight not only the direct electricity consumption of AI workloads but also the upstream and embodied emissions associated with providing AI services. 
 These works focus on the effect of AI models at the micro-level---at the level of the data center. But much is still unknown about the effect of AI power usage at a more macro-level such as on the power quality and power demand. 
 
 While research throughout the literature has attempted to analyze AI's impacts on the grid, referenced prior studies are unable to analyze the issue at the system level instead focusing on the statements that build up from computational requirements of specific models. We are able to analyze impacts at the grid level by utilizing market data and treating publication dates as treatment. As such, we propose an approach to analyze the demand impacts of AI models without information on the specific activities or makeup of any data centers.
 %Collectively, these contributions underscore that efficiency gains in model training and inference must be paired with structural shifts in energy procurement if meaningful emission reductions are to be achieved.
%Recently,the Energy and AI world energy outlook special report \cite{RN1} has highlighted the importance of coordinated energy planning to avoid supply bottlenecks at the grid. The goal of this paper is to address this gap by studying the impact of AI models beyond a single data center, and instead look at its effect on the power quality and power demands at the grids.
%while \cite{RN7} argue that accounting precision is a necessary condition for aligning AI deployment with broader sustainability objectives. This strand of the literature therefore links corporate procurement strategies with systemic issues of grid reliability, carbon intensity, and regulatory governance.